%==============================================================================
%  Chapter 2 Guide: How to Write a Literature Review
%  Compile: make  (or pdflatex chapter2-guide.tex twice)
%==============================================================================
\documentclass[11pt,a4paper]{article}

\usepackage[utf8]{inputenc}
\usepackage[T1]{fontenc}
\usepackage[a4paper,margin=1in]{geometry}
\usepackage{amsmath,amssymb}
\usepackage{xcolor}
\usepackage[most]{tcolorbox}
\usepackage{enumitem}
\usepackage{booktabs}
\usepackage{longtable}
\usepackage{tabularx}
\usepackage{fancyhdr}
\usepackage{titlesec}
\usepackage{listings}
\usepackage{array}
\usepackage[hidelinks]{hyperref}

\definecolor{forest}{HTML}{166534}
\definecolor{vibrant}{HTML}{22C55E}
\definecolor{deepforest}{HTML}{064E3B}
\definecolor{slate}{HTML}{1E293B}
\definecolor{codebg}{HTML}{F8FAFC}

\hypersetup{linkcolor=forest, urlcolor=forest}

\titleformat{\section}{\Large\bfseries\color{forest}}{\thesection}{0.6em}{}
\titleformat{\subsection}{\large\bfseries\color{deepforest}}{\thesubsection}{0.6em}{}
\titleformat{\subsubsection}{\normalsize\bfseries\color{slate}}{\thesubsubsection}{0.5em}{}

\setlist{leftmargin=1.4em, topsep=4pt, itemsep=2pt}

\lstset{
  basicstyle=\ttfamily\small,
  backgroundcolor=\color{codebg},
  frame=single,
  framerule=0.3pt,
  rulecolor=\color{slate!30},
  breaklines=true,
  columns=fullflexible,
  keepspaces=true,
  xleftmargin=4pt,
  xrightmargin=4pt,
  aboveskip=6pt,
  belowskip=6pt,
}

\newtcolorbox{tipbox}{
  enhanced, breakable, colback=forest!6, colframe=forest,
  boxrule=0.7pt, arc=2pt, left=8pt, right=8pt, top=6pt, bottom=6pt,
  fonttitle=\bfseries\color{white}, coltitle=white,
  title={Tip}, attach title to upper={\par\smallskip}}

\newtcolorbox{warnbox}{
  enhanced, breakable, colback=vibrant!8, colframe=vibrant!70!black,
  boxrule=0.7pt, arc=2pt, left=8pt, right=8pt, top=6pt, bottom=6pt,
  fonttitle=\bfseries, coltitle=deepforest,
  title={Watch out}, attach title to upper={\par\smallskip}}

\newtcolorbox{examplebox}{
  enhanced, breakable, colback=slate!4, colframe=slate!55,
  boxrule=0.6pt, arc=2pt, left=8pt, right=8pt, top=6pt, bottom=6pt,
  fonttitle=\bfseries, coltitle=slate,
  title={Example}, attach title to upper={\par\smallskip}}

\pagestyle{fancy}
\fancyhf{}
\fancyhead[L]{\small\textcolor{forest}{Chapter 2 Guide: Literature Review}}
\fancyhead[R]{\small\textcolor{forest}{\thepage}}
\renewcommand{\headrulewidth}{0.4pt}

\title{\textbf{Chapter 2 Guide}\\[0.35em]
\large How to Write a Literature Review\\[0.25em]
\normalsize A practical handbook for thesis and project writing}
\author{}
\date{September 2026}

\begin{document}
\maketitle

\begin{abstract}
\noindent
A literature review is a critical synthesis that positions your study, not a list of paper summaries.
This handbook follows established research-writing practice and walks through search, reading, thematic organisation, synthesis, gap writing, and revision.
Section~\ref{sec:worked} gives a domain-agnostic mini outline you can fill for any thesis topic; it is a \emph{template only} and does not replace a thesis Chapter~2.
\end{abstract}

\tableofcontents
\newpage

\section{What a literature review is (and is not)}

\subsection{What it is}
A literature review is a \textbf{structured argument} built from published work. You:
\begin{itemize}
  \item map what is already known on your topic;
  \item group findings into \textbf{themes};
  \item compare methods, settings, and limits;
  \item show where evidence is strong, weak, or missing;
  \item end with a \textbf{research gap} that your study addresses.
\end{itemize}
In research-writing terms (Hart; Ridley; and similar guides), the review \textbf{evaluates and organises} knowledge so the reader sees why your question is worth asking.

\subsection{What it is not}
\begin{center}
\begin{tabular}{@{}p{0.32\textwidth}p{0.58\textwidth}@{}}
\toprule
\textbf{Not this} & \textbf{Why it fails} \\
\midrule
Annotated bibliography & Lists papers one by one with no argument \\
Paper dump & ``Author A said\ldots\ Author B said\ldots'' with no link \\
Background story only & Context without critique or gap \\
Copy of abstracts & Summary without your analytical voice \\
Proof that you read a lot & Volume without selection or synthesis \\
\bottomrule
\end{tabular}
\end{center}

\begin{tipbox}
\textbf{Test:} if you remove the citations, does a clear story about the field still remain? If not, you are listing, not reviewing.
\end{tipbox}

%------------------------------------------------------------------------------
\section{Purpose in a thesis or project}

In a final-year or postgraduate thesis, the literature review usually does five jobs:
\begin{enumerate}
  \item \textbf{Situate the problem} --- show the issue exists and matters.
  \item \textbf{Define key terms} --- triage, acuity, pathway, embedding, Softmax, and so on.
  \item \textbf{Survey prior approaches} --- protocols, scores, NLP models.
  \item \textbf{Justify your method} --- why this model, colour map, or gate.
  \item \textbf{Lead into objectives} --- the gap must match what you will do next.
\end{enumerate}

\subsection{The funnel}
Keep this shape when outlining:
\begin{lstlisting}
Broad context (hospital pressure, triage need)
        |
Specific strands (vitals scores; NLP; domain models)
        |
Limits and mismatches (setting, data, missing local routing)
        |
Research gap
        |
Bridge to methodology / objectives
\end{lstlisting}
Presentation literature slides and thesis Chapter~2 should follow the same funnel: context \(\to\) strands \(\to\) gap.

%------------------------------------------------------------------------------
\section{Types of review}

Be honest about which type you are writing.

\begin{center}
\begin{tabular}{@{}p{0.22\textwidth}p{0.38\textwidth}p{0.28\textwidth}@{}}
\toprule
\textbf{Type} & \textbf{What it is} & \textbf{Typical for} \\
\midrule
Narrative (traditional) &
Thematic synthesis guided by your question; search structured but not always fully reproducible &
Most BSc / MPhil chapters \\
Systematic &
Pre-defined question, strict inclusion, exhaustive search, often PRISMA &
Evidence reviews; rare as full BSc chapter \\
Scoping &
Maps breadth of a field; less depth of appraisal &
Early exploration \\
Integrative &
Combines empirical and theoretical work into a framework &
Advanced conceptual chapters \\
\bottomrule
\end{tabular}
\end{center}

\begin{warnbox}
For most mathematics / applied AI final-year projects, write a \textbf{narrative thematic review}.
Do not call the chapter ``systematic'' unless you followed systematic methods end to end.
\end{warnbox}

%------------------------------------------------------------------------------
\section{Before writing: search and selection}
\label{sec:search}

Writing starts with \textbf{finding and filtering}, not with a blank chapter file.

\subsection{Clarify the review question}
Turn the topic into searchable sub-questions, for example:
\begin{itemize}
  \item How has free-text chief-complaint NLP been used for hospital triage?
  \item How do SATS and TEWS structure urgency and routing?
  \item What do Transformer / BioBERT models contribute to clinical text classification?
\end{itemize}

\subsection{Build keyword sets}
\begin{center}
\begin{tabular}{@{}p{0.28\textwidth}p{0.34\textwidth}p{0.28\textwidth}@{}}
\toprule
\textbf{Concept A} & \textbf{Concept B} & \textbf{Concept C} \\
\midrule
triage, acuity, SATS, TEWS &
chief complaint, NLP, text classification &
BioBERT, BERT, attention, LMIC, Ghana, pathway \\
\bottomrule
\end{tabular}
\end{center}
Combine with AND / OR. Example query:
\begin{lstlisting}
("chief complaint" OR "presenting complaint")
AND (triage OR acuity)
AND (NLP OR BERT OR BioBERT)
\end{lstlisting}

\subsection{Where to search}
\begin{itemize}
  \item Google Scholar
  \item PubMed / PMC (clinical and BioBERT lineage)
  \item IEEE Xplore / ACM (systems and NLP)
  \item University library databases
  \item Official manuals (e.g.\ SATS training manual) as primary protocol sources
\end{itemize}

\subsection{Inclusion and exclusion}
Write the rules down. Example:
\begin{itemize}
  \item \textbf{Include:} English (or working language); triage, pathways, or clinical NLP; peer-reviewed or recognised protocol; relevant to text or vitals at intake.
  \item \textbf{Exclude:} pure imaging AI with no text path; opinion pieces with no method; duplicates; topics that never touch your gap.
\end{itemize}

\subsection{Snowballing}
From 3--5 seed papers (e.g.\ BioBERT; a triage NLP review; SATS manual; a local overcrowding study):
\begin{itemize}
  \item check \textbf{references} (backward);
  \item check \textbf{who cited them} (forward on Scholar).
\end{itemize}

\subsection{Search log}
Keep a table:
\begin{center}
\begin{tabular}{@{}llllll@{}}
\toprule
Date & Database & Query & Hits & Kept & Notes \\
\midrule
\ldots & Scholar & chief complaint NLP triage & 120 & 8 & review useful \\
\bottomrule
\end{tabular}
\end{center}

\subsection{Stop rule}
Stop searching when new papers \textbf{repeat the same themes} and no longer change your gap statement.

%------------------------------------------------------------------------------
\section{Reading workflow}

Do not read every paper at the same depth.

\subsection{Three passes}
\begin{enumerate}
  \item \textbf{Skim} --- title, abstract, conclusion, figures. Keep / drop / maybe.
  \item \textbf{Structure} --- introduction, method, results, limits. Fill a note card.
  \item \textbf{Critique} --- what would fail if we applied this to \emph{your} setting?
\end{enumerate}

\subsection{Note card (one paper = one card)}
\begin{lstlisting}
Citation:
Year / venue:
Setting (country, hospital type):
Research question:
Method (data, model, vitals vs text):
Main finding (1-2 sentences):
Limitations (authors' + yours):
How I will use this (theme + gap role):
Quote worth paraphrasing (page):
\end{lstlisting}

\subsection{File notes by theme, not by author}
Use tags such as: overcrowding / SATS; TEWS / vitals; attention / BERT / BioBERT; NLP triage; pathways / routing.
When you write, open a \textbf{theme}, not an author alphabet.

%------------------------------------------------------------------------------
\section{Organisation patterns}

\begin{center}
\begin{tabular}{@{}p{0.22\textwidth}p{0.38\textwidth}p{0.28\textwidth}@{}}
\toprule
\textbf{Pattern} & \textbf{How it works} & \textbf{When} \\
\midrule
Thematic (preferred) &
Sections = ideas &
Almost always for thesis chapters \\
Chronological &
Early work \(\to\) recent work &
History of one narrow idea \\
Methodological &
Rule-based vs classical ML vs deep learning &
Comparing technique families \\
\bottomrule
\end{tabular}
\end{center}

\subsection{Why theme beats author-by-author}

\begin{examplebox}
\textbf{Weak:} Smith (2018) found X. Jones (2019) found Y. Lee (2020) proposed BioBERT.

\textbf{Stronger:} Clinical text models moved from bag-of-words classifiers to contextual encoders.
BioBERT continues BERT pre-training on PubMed/PMC so biomedical terms cluster more usefully in embedding space (Lee et al., 2020).
Triage applications of deep attention to chief complaints exist, but most use high-income corpora (Stewart et al., 2023).
\end{examplebox}

%------------------------------------------------------------------------------
\section{Summary vs synthesis}

\subsection{Summary (necessary but not enough)}
Restates what one source says:
\begin{quote}
Lee et al.\ (2020) introduced BioBERT by continuing BERT training on PubMed and PMC.
\end{quote}

\subsection{Synthesis (what markers look for)}
Connects sources, contrasts them, and points toward your problem:
\begin{quote}
Vital-sign scores such as TEWS give an objective colour band once measurements exist, but they do not read free text (SATS manual).
In parallel, NLP triage research classifies chief complaints with deep models, largely on high-income datasets (Stewart et al., 2023).
What remains under-specified for Ghanaian intake is an examinable chain from English complaint text to SATS colour and local pathway routing.
\end{quote}

\begin{center}
\begin{tabular}{@{}p{0.42\textwidth}p{0.48\textwidth}@{}}
\toprule
\textbf{Summary only} & \textbf{Synthesis} \\
\midrule
Lists findings & Relates findings to each other \\
Neutral report & Takes an analytical position \\
Ends anywhere & Ends at a gap or tension \\
``X said'' & ``X and Y together imply\ldots, but Z is missing'' \\
\bottomrule
\end{tabular}
\end{center}

\begin{tipbox}
\textbf{Rule of thumb:} each paragraph should answer ``so what for my project?'' in the last sentence or two.
\end{tipbox}

%------------------------------------------------------------------------------
\section{Critical appraisal}

For every important source, ask:
\begin{enumerate}
  \item \textbf{Setting} --- high-income ED vs LMIC teaching hospital; can results transfer?
  \item \textbf{Data} --- size, language, label source (nurse acuity vs diagnosis)?
  \item \textbf{Input} --- text only, vitals only, or multimodal?
  \item \textbf{Output} --- diagnosis, disposition, acuity, colour, pathway?
  \item \textbf{Evaluation} --- accuracy alone, or also recall on rare urgent classes, latency, safety?
  \item \textbf{Limits} --- what authors admit, and what they ignore.
  \item \textbf{Relevance to your gap} --- supports method, problem, or shows what is still missing?
\end{enumerate}

Prefer calm critique:
\begin{quote}
Most published NLP triage systems target high-income corpora; local department routing is rarely formalised.
\end{quote}
Avoid dismissive language such as ``these papers are useless.''

%------------------------------------------------------------------------------
\section{Chapter and section architecture}

\subsection{Typical thesis chapter shape}
\begin{enumerate}
  \item Short opener: what this chapter covers.
  \item Theme 1: broadest clinical or operational context.
  \item Theme 2: established hospital standard (e.g.\ vitals score).
  \item Theme 3: technical lineage (attention \(\to\) BERT \(\to\) domain model).
  \item Theme 4: applied NLP triage literature.
  \item Research gap: explicit, matching objectives.
  \item Bridge sentence to methodology.
\end{enumerate}

Prefer \textbf{few focused sections} over many tiny ones. Each section should carry one claim family.

\subsection{Presentation slides}
On a slide, show only the compressed funnel: 3--5 bullets plus the gap.
The chapter holds the synthesis; the slide holds the map.

%------------------------------------------------------------------------------
\section{Academic voice and paragraph craft}

\subsection{Paragraph skeleton}
\begin{enumerate}
  \item \textbf{Topic sentence} --- your claim about the literature.
  \item \textbf{Evidence} --- 1--3 citations with brief findings.
  \item \textbf{Analysis} --- compare, limit, or interpret.
  \item \textbf{So-what} --- link to problem, method, or gap.
\end{enumerate}

\begin{examplebox}
\textbf{Topic:} Text arrives before vitals at intake.\\
\textbf{Evidence:} Chief complaints are unstructured prose; TEWS needs measured parameters.\\
\textbf{Analysis:} Vitals-only pathways cannot cover the first minutes of arrival.\\
\textbf{So-what:} A text path is needed alongside TEWS, not instead of it.
\end{examplebox}

\subsection{Tense}
\begin{center}
\begin{tabular}{@{}lp{0.62\textwidth}@{}}
\toprule
\textbf{Use} & \textbf{For} \\
\midrule
Past & Specific completed studies (``Lee et al.\ (2020) trained\ldots'') \\
Present & Established knowledge (``SATS uses five colour bands\ldots'') \\
Present perfect & Line of research up to now (``Researchers have applied NLP to triage\ldots'') \\
\bottomrule
\end{tabular}
\end{center}

\subsection{Hedging and signposting}
Use precise hedges: suggests, reports, associated with, under-specified, limited evidence.\\
Avoid overclaiming: proves, all hospitals, always.\\
Signpost: ``A related limitation is\ldots''; ``In contrast to vitals-based scores\ldots''; ``These strands motivate the gap below.''

Avoid filler openings (``In this modern era of AI\ldots'') and habitual em dashes; prefer commas, colons, or new sentences.

%------------------------------------------------------------------------------
\section{Citations and ethics}

\begin{itemize}
  \item Cite to support a claim; definitions from manuals need the manual.
  \item \textbf{Default:} paraphrase and cite. Quote rarely (precise protocol wording).
  \item Do not copy sentence structure with only synonym swaps.
  \item Citation clusters are fine only if sources share the same point.
  \item Prefer primary sources; if citing via a review, follow department style.
  \item Do not invent citations; keep PDFs or stable links in a reference manager.
\end{itemize}

%------------------------------------------------------------------------------
\section{Writing the research gap}
\label{sec:gap}

The gap is the payoff of the review. Objectives and method must answer it.

\subsection{Gap formula}
\begin{lstlisting}
Prior work establishes A and B.
However, C remains missing / under-specified for setting S.
Therefore this study does D.
\end{lstlisting}

\subsection{Weak vs strong}

\begin{examplebox}
\textbf{Weak:} Little research has been done on hospital chatbots in Ghana.\\
(Vague; may be false; does not specify what is missing.)

\textbf{Stronger:} Existing work separates vital-sign scores (TEWS), manual SATS colour assignment, and NLP classifiers trained mainly on high-income corpora.
There is no unified, examinable account of how free-text chief complaints at a Ghanaian teaching-hospital-style intake become a SATS colour and a local pathway card, while TEWS remains the nurse-facing vital standard.
\end{examplebox}

\subsection{Match gap to pipeline (illustrative)}
If your contribution is an examinable chain such as
\[
X \;\to\; \hat{\mathbf{y}} \;\to\; C_{\mathrm{NLP}} \;\to\; P(C),
\]
then the gap paragraph should mention \textbf{text \(\to\) colour \(\to\) local routing}, not a vague ``AI in hospitals.''

Aim for \textbf{one coherent gap}, even if several missing pieces feed into it.

%------------------------------------------------------------------------------
\section{Common mistakes checklist}

\begin{itemize}[label=$\square$]
  \item Author-by-author laundry list
  \item No critique (only praise or neutral summary)
  \item Themes missing; chronology used with no reason
  \item Gap too broad (``no one studied AI'')
  \item Gap does not match stated objectives
  \item Method not justified by literature
  \item Local / LMIC transfer ignored when claiming hospital relevance
  \item Mixing TEWS and NLP without clarifying different inputs
  \item Orphan citations (text vs bibliography mismatch)
  \item Over-quoting
  \item Informal tone or unsupported superlatives
  \item Sections that never return to the research question
  \item Ending without an explicit gap subsection
\end{itemize}

%------------------------------------------------------------------------------
\section{Revision checklist}

\subsection{Structure}
\begin{itemize}[label=$\square$]
  \item Opener states purpose of the chapter
  \item Sections are thematic and in funnel order
  \item Each section ends with a clear takeaway
  \item Gap subsection is explicit and specific
  \item Bridge to methodology exists
\end{itemize}

\subsection{Argument}
\begin{itemize}[label=$\square$]
  \item Synthesis appears in most paragraphs
  \item Contradictions or limits are acknowledged
  \item Gap follows logically from the themes
\end{itemize}

\subsection{Scholarship}
\begin{itemize}[label=$\square$]
  \item Core seed papers present and correctly cited
  \item Protocol manuals cited where colours / timers are defined
  \item Recent and foundational sources both appear where needed
  \item Bibliography matches department style
\end{itemize}

\subsection{Clarity and alignment}
\begin{itemize}[label=$\square$]
  \item Terms defined on first use; abbreviations expanded once
  \item Sentences readable aloud
  \item Problem, objectives, and gap use the same vocabulary
  \item Presentation literature slide does not contradict the chapter
\end{itemize}

%------------------------------------------------------------------------------
\section{Mini worked outline (illustrative template)}
\label{sec:worked}

\begin{warnbox}
Use this as a \textbf{domain-agnostic shape}. Replace every bracketed placeholder with your own topic, sources, and wording.
It is a \textbf{template}, not a substitute for thesis Chapter~2.
\end{warnbox}

\paragraph{Theme A --- Broader context / problem setting.}
Claim: [why the topic matters in the real world or in the science].\\
Evidence types: [background studies; policy; operational or scientific context papers].\\
So-what: [what this implies for your research question].

\paragraph{Theme B --- Established standard or classical approach.}
Claim: [the accepted method, protocol, baseline, or classical theory in the field].\\
Evidence types: [manuals, standard algorithms, widely cited baselines].\\
Limit: [what this standard does \emph{not} cover that your project needs].\\
So-what: [why a further strand of literature is required].

\paragraph{Theme C --- Core technical / theoretical lineage.}
Claim: [the key models, theorems, or frameworks your method builds on].\\
Evidence types: [foundational papers; textbooks; method surveys].\\
So-what: [why this lineage justifies the tools you will use].

\paragraph{Theme D --- Applied / closely related studies.}
Claim: [recent work closest to your exact task].\\
Evidence types: [empirical papers; applied case studies; reviews of near neighbours].\\
Limit: [setting, data, language, or scope that does not match yours].\\
So-what: [near work is not the same as solving your specific problem].

\paragraph{Theme E --- Research gap (payoff).}
Prior: Themes A--D establish [A] and [B].\\
Missing: [X] remains under-specified for setting [S].\\
Bridge: this study addresses that gap by [one sentence on your contribution / method chain].

\begin{tipbox}
\textbf{Viva closer:} ``The literature gives established approaches and related models; what was missing for us was a clear, examinable account of [your specific chain or question].''
\end{tipbox}

%------------------------------------------------------------------------------
\section{Quick templates (copy and fill)}

\subsection{Search log row}
\begin{lstlisting}
Date:
Database:
Query:
Results (n):
Titles screened:
Full texts read:
Kept (citations):
Reason excluded (examples):
\end{lstlisting}

\subsection{Paper note card}
\begin{lstlisting}
Citation:
Setting:
Input (text / vitals / both):
Output (acuity / colour / diagnosis / other):
Method:
Finding:
Limit:
Use in my review (theme):
Supports gap? (yes/no + how):
\end{lstlisting}

\subsection{Synthesis paragraph}
\begin{lstlisting}
[Topic sentence: claim about the field.]
[Source 1 finding + citation.]
[Source 2 finding + citation; agreement or contrast.]
[Limit or transfer issue for my setting.]
[So-what: link to problem, method, or gap.]
\end{lstlisting}

\subsection{Gap paragraph}
\begin{lstlisting}
Prior work establishes [A] and [B] (citations).
These approaches leave [C] under-specified for [setting].
In particular, [what exact chain or artefact is missing].
This study addresses that gap by [one sentence on contribution].
\end{lstlisting}

\subsection{Theme section outline}
\begin{lstlisting}
Section title:
One-sentence section claim:
3-6 sources to weave:
Limits to mention:
Closing so-what sentence:
\end{lstlisting}

%------------------------------------------------------------------------------
\section{Further reading (craft guides)}

These are guides on \emph{how} to write reviews (consult your library for editions):
\begin{itemize}
  \item Chris Hart --- \emph{Doing a Literature Review}
  \item Diana Ridley --- \emph{The Literature Review: A Step-by-Step Guide for Students}
  \item Jill Jesson et al.\ --- \emph{Doing Your Literature Review}
  \item PRISMA materials --- only if you attempt a systematic-style review
\end{itemize}
Topic citations (BioBERT, SATS, Stewart, and so on) belong in the thesis bibliography, not necessarily here.

\section*{Final reminder}
A good literature review is judged by \textbf{clarity of argument and honesty of gap}, not by how many papers you name.
Select carefully, synthesise by theme, critique transfer to your setting, and make the gap match your objectives.

\end{document}
